% ============================================
% Chapter 2: Literature Review
% ============================================

\chapter{Literature Review}
\label{ch:literature-review}



Literature Review

\section{Review Strategy}


The literature review was restricted to sources that directly contribute to GreenLab's core
design concerns: IoT energy management, occupancy-driven smart-building control, YOLObased visual detection, ESP32 or Arduino-class device control, laboratory or classroom automation, and web dashboards for real-time IoT monitoring. Sources were not included
merely because they used IoT terminology or addressed building energy in general. Each
considered source is preceded by an explicit relevance analysis comment, and only relevant
papers and projects are reviewed in detail. The review strategy follows an inclusion and exclusion methodology: a source is included if it addresses at least two of GreenLab's design
concerns and provides verifiable technical content that can be accurately summarised and
attributed; a source is excluded if it addresses only tangentially related topics, if its full text
is unavailable for verification, or if more recent and directly relevant sources supersede it.

Eight sources met this standard and are reviewed in Sections 2.2.1 through 2.2.8.
Each review follows the same analytical structure: a description of the methodology and
the hardware or software choices made by the original authors, the core findings reported in
the source, an explicit breakdown of advantages and disadvantages relative to GreenLab's
design concerns, and a statement of what GreenLab adopts, adapts, or deliberately departs
from.

Section 2.3 consolidates all eight sources into a single comparative matrix, and Section 2.4 synthesises the literature gap that GreenLab is designed to close.




\section{Related Work}



11



\subsection{Challa et al. (2025) --- Multimodal Classroom Occupancy for HVAC}


Methodology and Technical Summary Challa et al. proposed a classroom occupancydetection method that combines real-time video analysis with gas-sensor data for HVAC
optimisation [?]. The paper is directly relevant to GreenLab because both systems treat
educational spaces as variable-occupancy environments and use visual occupancy information to support energy-related control decisions. The study integrates a webcam for visual
capture, an MS1100 gas sensor for environmental monitoring, Arduino-class hardware for
data collection, Python and OpenCV for image processing, and a YOLOv4 detection model
for person counting. The authors report that the model was trained using more than 20,000
labelled human face images and that multimodal sensing was used to improve reliability
under conditions where visual detection alone may be affected by lighting, occlusion, or
camera blind spots. The gas sensor provides an independent signal that can corroborate or
contradict the camera-derived occupancy estimate, reducing the risk that a single modality
failure leads to an incorrect HVAC response [?].

The methodology is important for GreenLab because it demonstrates the value of combining environmental signals with computer vision. GreenLab similarly avoids reliance on
one sensing method by using environmental sensors through the ESP32 and camera frames
through the AI processing service. The gas-sensor approach used by Challa et al. is analogous to GreenLab's use of the DHT11 temperature sensor and the HC-SR04 ultrasonic
sensor as corroborating inputs alongside the YOLO-derived occupancy count. The key insight from this paper is that HVAC or ventilation decisions should not be based solely on
camera detection, because camera accuracy can degrade under variable lighting, occlusion
from furniture or equipment, and camera placement constraints that are common in laboratory environments.

Challa et al. also emphasise that occupancy information can be connected to dynamic
HVAC control, which supports GreenLab's decision to link occupancy and temperature
context to fan or ventilation states. The results reported in the cited paper, including improved HVAC response times and reduced false-positive occupancy classifications, are not
transferred to GreenLab as project results; rather, they justify the design principle that multimodal occupancy estimation can be more robust than isolated motion sensing. GreenLab
extends the multimodal principle by adding natural-light sensing for lighting control, a

12
GreenLab: Intelligent Laboratory Energy Management System


backend rule engine for governance, and a role-based dashboard for human supervision,
none of which are present in the Challa et al. system.



\subsubsection*{Advantages}



\begin{itemize}
  \item Multimodal corroboration. Combining gas sensing with visual detection reduces
the risk that a single sensor failure (poor lighting, occlusion) produces an incorrect
HVAC decision.
\end{itemize}

\begin{itemize}
  \item Large training dataset. The YOLOv4 model was trained on more than 20,000 la-
belled images, giving the detector a substantial data foundation relative to many com-
parable classroom studies.
\end{itemize}

\begin{itemize}
  \item Direct HVAC linkage. Occupancy output is connected explicitly to dynamic HVAC
control rather than left as a standalone detection result, demonstrating a complete
sense-to-actuate pipeline.
\end{itemize}



\subsubsection*{Disadvantages / Limitations}



\begin{itemize}
  \item Privacy-sensitive training data. Reliance on face-based image detection (rather
than whole-person or silhouette detection) raises privacy considerations that are less
acceptable in many educational deployment contexts.
\end{itemize}

\begin{itemize}
  \item Gas sensor suitability. The MS1100 gas sensor may not be appropriate or neces-
sary for all educational environments, adding hardware cost and calibration overhead
without a clearly transferable benefit to a laboratory setting.
\end{itemize}

\begin{itemize}
  \item No governance layer. The system has no web-based dashboard, no role-based access
control, and no audit trail for override or administrative action --- it is a detection-
and-actuation pipeline, not an operationally governed platform.
\end{itemize}



\subsubsection*{Relevance to GreenLab}

departs from Challa et al.'s specific sensor choices: person detection is used instead of face
detection to reduce privacy sensitivity, and temperature and light sensors are used instead
of a gas sensor, since they are less privacy-sensitive and more directly tied to GreenLab's

13


lighting and ventilation decisions. Most significantly, GreenLab adds the complete backend
and dashboard governance layer that this paper does not attempt.



\subsection{Chaudhari et al. (2024) --- Occupancy Detection Survey for Smart}

Buildings

Methodology and Technical Summary Chaudhari, Xiao, Cheng, and Li presented a
broad survey of occupancy detection technologies for smart buildings using IoT sensors
[?]. The survey is relevant because GreenLab requires reliable occupancy information for
lighting, ventilation, monitoring, and safety-related decisions, and the survey provides a
comprehensive taxonomy of the available sensing methods, their strengths, and their limitations. The paper explains that occupancy information can reduce energy consumption by
automating lighting and HVAC systems, improve indoor air quality by operating ventilation
only where needed, and support security and emergency response by providing real-time
occupancy data. The survey covers five principal sensing methods: passive infrared sensors, carbon dioxide sensors, camera-based detection, Wi-Fi-based device counting, and
Bluetooth Low Energy proximity sensing [?].

The survey identifies several practical limitations that are directly relevant to GreenLab's design. Sensor placement affects detection accuracy because a PIR sensor positioned
near an air conditioner vent may produce false positives from moving air, while a camera
placed behind a projector screen may have an obstructed field of view. Privacy concerns
limit the deployment of camera-based systems in some educational contexts, and real-time
constraints require that occupancy decisions be made within the control-loop budget rather
than after lengthy batch processing. The survey also highlights interference between sensing modalities, such as Wi-Fi counting that may overestimate occupancy when students
carry multiple wireless devices.

The survey's contribution to GreenLab is conceptual and architectural rather than implementational, since the paper does not describe a deployed system. It supports the decision to treat occupancy detection as a system-level capability rather than as a single sensor
reading. If occupancy is a system-level capability, then the system must be designed to fuse
evidence from multiple sensors, resolve contradictions, and present a consistent occupancy
state to downstream control logic. It also reinforces the need for privacy-aware design,

14
GreenLab: Intelligent Laboratory Energy Management System


maintainability, and fault handling. GreenLab therefore uses a backend decision layer and
dashboard audit trail so that sensor readings, vision-based counts, manual overrides, and
faults can be interpreted in context rather than in isolation. The survey provides the empirical and theoretical foundation for GreenLab's multimodal approach, while GreenLab
provides the governance layer that the survey identifies as a practical gap in many existing
systems.




\subsubsection*{Advantages}



\begin{itemize}
  \item Comprehensive taxonomy. Covers five distinct sensing modalities with a consistent
comparative framework, giving a structured basis for evaluating sensor trade-offs
rather than relying on a single technology's marketing claims.
\end{itemize}


\begin{itemize}
  \item Practical limitation analysis. Documents concrete, deployment-relevant failure
modes (sensor placement, occlusion, multi-device overcounting) rather than only the-
oretical accuracy figures.
\end{itemize}


\begin{itemize}
  \item System-level framing. Treats occupancy detection as a property of the system as
a whole rather than of any one sensor, supporting fusion-based architectures like
GreenLab's.
\end{itemize}




\subsubsection*{Disadvantages / Limitations}



\begin{itemize}
  \item No deployed system. As a survey, the paper does not describe or validate an end-
to-end automation platform, so its conclusions are necessarily general rather than
specific to a laboratory deployment.
\end{itemize}


\begin{itemize}
  \item No governance discussion. The survey focuses on detection accuracy and does not
address authentication, role-based access, override management, or audit logging.
\end{itemize}


\begin{itemize}
  \item No quantitative system comparison. Because the survey aggregates findings from
many underlying studies, it does not provide a single, consistent benchmark across
all five modalities under identical conditions.
\end{itemize}

15



\subsubsection*{Relevance to GreenLab}

it is excluded from the system comparison table in Section 2.3 but is drawn on throughout
this chapter to justify GreenLab's multimodal sensing rationale and its emphasis on fusing
rather than relying on a single sensing modality.





\textit{Figure 2.1: Occupancy-detection taxonomy and sensing modalities discussed by Chaudhari et al.}

for smart-building applications.




\subsection{Zhang et al. (2026) --- Modified Lightweight YOLOv8 for Indoor}

Occupancy

Methodology and Technical Summary Zhang et al. proposed a modified lightweight
YOLOv8 model for fast and precise indoor occupancy detection [?]. The paper is relevant
to GreenLab's AI processing layer because GreenLab depends on person detection or occupancy estimation from laboratory camera frames. The authors identify a common problem
in image-based indoor occupancy detection: small human targets and resource constraints
can reduce real-time performance. In a typical laboratory or classroom, occupants may

16
GreenLab: Intelligent Laboratory Energy Management System


be seated at desks with only their upper bodies visible, they may be partially occluded by
monitors and equipment, and the camera may be mounted at a height and angle that makes
person detection more challenging than in open pedestrian scenes. These factors reduce
the detection accuracy of standard YOLO models that are trained primarily on outdoor or
full-body datasets [?].

The proposed architecture modifies YOLOv8 using feature-map improvements and
convolutional adjustments, and the paper reports reduced parameters and computational
complexity relative to a baseline YOLOv8 model while improving detection performance
on the evaluated indoor dataset. The authors demonstrate that lightweight architectures
can achieve competitive detection accuracy with fewer computational resources, which is
important for deployment scenarios where a dedicated GPU may not be available or where
multiple camera streams must be processed concurrently. The paper also discusses the
trade-off between model complexity and inference speed, noting that real-time occupancy
detection for building automation requires inference times compatible with the control-loop
budget, typically a few seconds per frame.

For GreenLab, this paper supports the selection of a YOLO-family detector and informs future optimisation. The present report does not claim that GreenLab implements the
modified architecture exactly. Instead, the paper is used to justify why lightweight detection is a meaningful concern if the AI layer is deployed on constrained local hardware or
expected to process multiple laboratory streams. GreenLab's design therefore separates the
AI processing service from the ESP32 controller, allowing the model to be improved without changing the embedded relay-control firmware. This separation is consistent with the
layered architecture principle described in the project objectives and with the broader literature on modular IoT systems. The paper also supports a future optimisation direction where
GreenLab could adopt a lightweight YOLO variant to reduce inference latency and server
resource consumption, particularly as the number of monitored laboratories increases.




\subsubsection*{Advantages}



\begin{itemize}
  \item Reduced computational cost. Reports lower parameter count and computational
complexity relative to baseline YOLOv8 while improving detection performance on
indoor data, directly relevant to constrained-hardware deployment.
\end{itemize}

17


\begin{itemize}
  \item Indoor-specific optimisation. Explicitly targets the small-target and occlusion prob-
lems that arise in seated, desk-based occupancy scenarios --- the same scenario class
as a computer science laboratory.
\end{itemize}



\begin{itemize}
  \item Compatibility with real-time control-loop budgets. Demonstrates inference times
suitable for building-automation control loops rather than only offline batch analysis.
\end{itemize}





\subsubsection*{Disadvantages / Limitations}




\begin{itemize}
  \item Algorithmic paper, not a deployed system. The paper validates a detection model
in isolation; it does not address sensor fusion, relay control, governance, or any end-
to-end automation workflow.
\end{itemize}



\begin{itemize}
  \item Single-dataset evaluation. Results are reported on the authors' evaluated indoor
dataset, and transferability to a specific laboratory's camera placement, lighting, and
occupant density is not independently verified.
\end{itemize}



\begin{itemize}
  \item No discussion of multimodal fusion. The model is evaluated purely on visual input,
without considering how its output should be reconciled with environmental sensor
signals.
\end{itemize}





\subsubsection*{Relevance to GreenLab}

complete automation platform. GreenLab does not claim to implement Zhang et al.'s exact
modified architecture; rather, the paper justifies why model lightness is a legitimate design concern and supports a stated future-work direction (Chapter 7) in which a lightweight
YOLO variant could be adopted to reduce inference latency as the number of monitored
laboratories scales.

18
GreenLab: Intelligent Laboratory Energy Management System





\textit{Figure 2.2: Modified lightweight YOLOv8 architecture used for indoor occupancy detection by}

Zhang et al.




\subsection{Waghchoure et al. (2025) --- IoT Classroom Appliance Control}

with PIR and Camera

Methodology and Technical Summary Waghchoure et al. described a classroom energysaving system using IoT, PIR motion sensing, camera-based detection, a microcontroller,
and relay control for lights and fans [?]. The paper is relevant because it addresses the
same educational-building problem that GreenLab targets: appliances may remain active
in unoccupied rooms, increasing electricity bills and environmental impact. The proposed
system detects human presence using a PIR sensor, checks whether devices are on or off,
controls appliances based on detection, and continuously monitors entry and exit activity.
The camera is activated when the PIR sensor detects motion, providing visual verification

19


that reduces false positives from non-human heat sources such as sunlight or heating vents
[?].

The methodology contributes to GreenLab by showing a practical classroom automation pattern: low-cost sensing triggers device control through relays. The PIR-and-camera
combination demonstrates that sensor-triggered camera activation can reduce continuous
camera processing overhead while still providing visual verification when it is most needed.
GreenLab adopts a similar approach in principle: the HC-SR04 ultrasonic sensor and other
embedded sensors provide the initial activity signal, and the YOLO service performs detailed occupancy counting when triggered. However, GreenLab extends this pattern significantly by adding a backend rule engine for governance, role-based user access, laboratory
assignment workflows, audit logging, and a web dashboard for instructors and administrators. The Waghchoure et al. system operates as a standalone microcontroller application
without a backend server, without user accounts, and without operational auditability.

The paper also reveals several practical challenges that GreenLab must address. PIR
sensors have limited detection range and may not detect stationary occupants, leading to
false vacancy detection. Camera-based verification introduces latency that must be managed within the control-loop budget. Relay control must handle the physical characteristics
of the controlled appliances, such as inrush current for fluorescent lighting and minimum
on-off cycling intervals for fans. GreenLab addresses these challenges through grace periods before vacancy-based appliance shutdown, confidence thresholds for YOLO occupancy
decisions, and configurable command intervals that prevent rapid relay cycling. These design choices are directly informed by the limitations identified in the Waghchoure et al.
work.




\subsubsection*{Advantages}



\begin{itemize}
  \item Low processing overhead. PIR-triggered camera activation avoids continuous video
processing, conserving compute resources until visual verification is actually needed.
\end{itemize}


\begin{itemize}
  \item Practical false-positive reduction. Camera verification after a PIR trigger filters
out non-human heat sources (sunlight, heating vents) that would otherwise cause
unnecessary appliance activation.
\end{itemize}

20
GreenLab: Intelligent Laboratory Energy Management System


\begin{itemize}
  \item Demonstrated low-cost feasibility. Confirms that commodity microcontroller and
relay hardware is sufficient for classroom-scale appliance automation.
\end{itemize}



\subsubsection*{Disadvantages / Limitations}



\begin{itemize}
  \item No backend server or persistence. The system operates entirely on the microcon-
troller, with no central data store, no historical record of occupancy or appliance state,
and no way to review past behaviour.
\end{itemize}

\begin{itemize}
  \item No user accounts or access control. Anyone with physical or network access to the
system can interact with it; there is no concept of an instructor or administrator role.
\end{itemize}

\begin{itemize}
  \item Stationary-occupant blind spot. Like all PIR-based systems, it cannot reliably de-
tect a seated, motionless occupant, risking false vacancy classification during quiet
study or exam periods.
\end{itemize}



\subsubsection*{Relevance to GreenLab}

implemented through the HC-SR04 sensor initiating a YOLO-based check, but replaces
the standalone-microcontroller model with a full backend and dashboard. The grace periods, confidence thresholds, and relay-cycling safeguards in GreenLab's design are direct,
traceable responses to the limitations this paper documents.



\subsection{Wibowo et al. (2023) --- IoT Energy Management in a University}

Laboratory

Methodology and Technical Summary Wibowo et al. presented an IoT-based energy
management system designed specifically for university laboratory environments [?]. The
paper is relevant to GreenLab because it targets the same environment, university laboratories, and the same goal, reducing energy waste through IoT-based monitoring and automation. The authors describe a centralised IoT sensor network that monitors laboratory
appliance status, environmental conditions, and energy consumption, and uses this data to
automate appliance switching and provide monitoring dashboards for laboratory administrators. The system demonstrates that centralised monitoring reduces energy waste in

21


academic labs by ensuring that appliances are switched off when laboratories are not in use
and that energy consumption patterns are visible to decision-makers [?].

The key contribution of this paper to GreenLab is the validation of the university laboratory context as a meaningful target for IoT energy management. While many IoT energy
systems target commercial offices or residential buildings, university laboratories have distinct characteristics: irregular schedules, variable occupancy, shared equipment, and multiple responsible parties. Wibowo et al. show that a centralised monitoring approach can
address these characteristics by providing real-time visibility and automated control.

However, the paper's system does not include computer-vision occupancy detection,
role-based access control, manual override with audit trails, or laboratory transfer workflows. These governance features are essential for a real faculty deployment where instructors must retain pedagogical authority, administrators must manage schedules and approvals, and all actions must be auditable.

GreenLab extends the Wibowo et al. approach by adding YOLO-based occupancy
detection that provides person counts rather than binary occupied or vacant signals, a rolebased dashboard that separates instructor and administrator functions, override management with time-limited expiry and audit logging, and transfer request workflows that prevent informal room changes from bypassing access control. The centralised monitoring
principle is retained, but the governance layer is significantly enhanced to meet the operational requirements of a faculty environment where accountability and safety are as important as energy savings.



\subsubsection*{Advantages}



\begin{itemize}
  \item Direct context match. Targets the university laboratory environment specifically,
rather than a generic office or residential setting, making its findings about irregular
schedules and shared equipment directly transferable to GreenLab.
\end{itemize}

\begin{itemize}
  \item Centralised visibility. Provides administrators with a single monitoring point for
appliance status and energy consumption across the laboratory network, rather than
per-room isolated control.
\end{itemize}

\begin{itemize}
  \item Demonstrated energy-waste reduction. Confirms that centralised IoT monitoring,
\end{itemize}

22
GreenLab: Intelligent Laboratory Energy Management System


even without vision-based occupancy detection, can reduce energy waste in an aca-
demic laboratory setting.



\subsubsection*{Disadvantages / Limitations}



\begin{itemize}
  \item Binary occupancy signal only. Without computer-vision integration, the system can
only report occupied/vacant states rather than person counts or occupancy classifica-
tion levels.
\end{itemize}

\begin{itemize}
  \item No role-based access control. All users of the monitoring dashboard appear to have
equivalent access, without separation between instructor and administrator privileges.
\end{itemize}

\begin{itemize}
  \item No override audit trail or transfer workflow. The paper does not describe a mech-
anism for authorised manual override with expiry, nor a workflow for handling labo-
ratory transfer requests.
\end{itemize}



\subsubsection*{Relevance to GreenLab}

sufficient, condition for a real faculty deployment --- the governance features absent from
this paper (role separation, override audit trails, transfer workflows) are precisely the features GreenLab adds on top of the validated centralised-monitoring pattern.



\subsection{Ahmed et al. (2020) --- IoT Classroom Automation with Sensor-}

Driven Control

Methodology and Technical Summary Ahmed et al. described an IoT-based classroom
automation system that uses motion detection, environmental sensors, a microcontroller,
and a web dashboard for classroom appliance switching [?]. The paper is relevant because
it demonstrates the feasibility of low-cost IoT classroom automation using off-the-shelf
components. The system detects classroom occupancy using motion sensors, reads environmental conditions such as temperature and light level, and uses a microcontroller to
switch lights and fans through relays. A web dashboard provides remote monitoring and
manual control, allowing instructors to check classroom status and override automatic decisions from a web browser [?].

23


The paper highlights several challenges that are directly relevant to GreenLab's design. Sensor placement affects detection reliability, because a motion sensor mounted in
a corner may not detect occupants near the door, and a temperature sensor placed near a
window may be affected by solar gain. False positives from non-human motion sources,
such as curtains moved by air conditioning, can trigger unnecessary appliance activation.
Delayed response after class changes, when the motion sensor has not yet detected vacancy,
can leave appliances running for extended periods after occupants have left. Ahmed et al.
note that these challenges can be partially mitigated through sensor calibration, debouncing, and grace periods, but that a more robust solution would require additional sensing
modalities and a centralised rule engine.

GreenLab addresses these challenges through multimodal sensing that combines ultrasonic, temperature, and light sensors with YOLO-based visual verification, grace periods that delay vacancy-based shutdown until occupancy absence is confirmed, confidence
thresholds that filter transient detections, and a backend rule engine that applies consistent,
auditable control logic. The web dashboard in the Ahmed et al. system provides basic
remote control but does not implement role-based access, time-limited override with automatic expiry, laboratory transfer approvals, or audit logging. GreenLab provides all of these
governance features, ensuring that the convenience of remote control does not compromise
operational accountability or safety.




\subsubsection*{Advantages}



\begin{itemize}
  \item Early validation of web-based remote control. Demonstrates, as early as 2020,
that a web dashboard can provide useful remote monitoring and manual override
for classroom automation, anticipating the dashboard-centred pattern GreenLab also
follows.
\end{itemize}

\begin{itemize}
  \item Transparent limitation analysis. The authors explicitly document sensor-placement
and false-positive challenges rather than presenting an idealised result, providing
concrete, citable design lessons.
\end{itemize}

\begin{itemize}
  \item Low-cost, off-the-shelf feasibility. Confirms that classroom automation does not
require specialised or expensive hardware to achieve basic occupancy-driven control.
\end{itemize}

24
GreenLab: Intelligent Laboratory Energy Management System



\subsubsection*{Disadvantages / Limitations}


\begin{itemize}
  \item No role-based access control. The web dashboard's manual override is available
without distinguishing between instructor and administrator authority.
\end{itemize}

\begin{itemize}
  \item No time-limited override or audit logging. Manual overrides, once issued, have no
documented expiry mechanism or persistent audit record.
\end{itemize}

\begin{itemize}
  \item Single-modality occupancy detection. Relies on motion sensing alone, without a
corroborating modality, leaving it exposed to the same false-positive and stationary-
occupant problems documented elsewhere in this chapter.
\end{itemize}



\subsubsection*{Relevance to GreenLab}

periods alone are insufficient mitigations for single-modality sensing; this is part of the
justification for combining ultrasonic and YOLO-based detection rather than relying on
motion sensing in isolation. The basic remote-override pattern in this paper is the conceptual precursor to GreenLab's fully governed override mechanism, with time-limited expiry,
role validation, and audit logging added on top.


\subsection{Ingole et al. (2024) --- ESP32 Classroom Automation}


Methodology and Technical Summary The IRJMETS classroom automation project
by Ingole et al. integrates ESP32 technology, LDR sensing, temperature sensing, relaycontrolled lighting and fan systems, entry management, two operational modes, and a
Wi-Fi manual control interface [?]. This project is highly relevant because its hardware
stack closely matches the GreenLab hardware concept. The ESP32 serves as the central
controller for sensor acquisition and actuator control, reading LDR values to determine
natural-light conditions, reading temperature values to determine ventilation requirements,
and activating relays to switch lights and fans. The system supports two operating modes:
an automatic mode in which the ESP32 applies sensor-based control rules independently,
and a manual mode in which the user can override automatic decisions through a Wi-Fi
interface [?].

The two-mode design is a useful reference for GreenLab because GreenLab also needs
to support both automatic and manual operation. However, the Ingole et al. implementation

25


treats manual control as a direct Wi-Fi command to the ESP32, without authentication, rolebased access, time-limited expiry, or audit logging. Any device on the same Wi-Fi network
can send a control command to the ESP32, which is acceptable for a prototype but not for a
production faculty system where unauthorised control could create safety hazards or operational disruptions. GreenLab adopts the embedded-control philosophy of this project but
places it inside a broader multi-tier architecture. In GreenLab, the ESP32 does not operate
as an isolated classroom controller; it communicates with an ASP.NET backend, receives
authorised control states, reports sensor data, and participates in audit-backed workflows.
Manual override is issued through the authenticated dashboard, validated by the backend,
recorded in the audit log, and subject to configurable expiry and cancellation policies.

The hardware similarity extends to the sensor types and relay configuration. Both
systems use an LDR for light sensing, a temperature sensor for ventilation decisions, and
a relay module for light and fan control. GreenLab adds the HC-SR04 ultrasonic sensor
for entry and activity detection, an IP camera for YOLO occupancy verification, and additional relay channels for multi-fan control. These additions are motivated by the limitations
of the Ingole et al. approach: without occupancy verification, the system cannot distinguish between a laboratory that is scheduled but empty and one that is actually occupied,
and without visual verification, sensor-based occupancy decisions are vulnerable to false
positives and false negatives.



\subsubsection*{Advantages}


\begin{itemize}
  \item Closest hardware match to GreenLab. Uses the same core controller (ESP32) and
overlapping sensor types (LDR, temperature), making this the most directly transfer-
able hardware reference in the literature reviewed.
\end{itemize}

\begin{itemize}
  \item Working dual-mode design. Demonstrates a functional automatic/manual mode
switch, validating the basic premise that an embedded controller can support both
autonomous and user-directed operation.
\end{itemize}

\begin{itemize}
  \item Simple, low-latency override path. Direct Wi-Fi commands to the ESP32 achieve
minimal latency for manual control, since no intermediate server is involved.
\end{itemize}



\subsubsection*{Disadvantages / Limitations}


26
GreenLab: Intelligent Laboratory Energy Management System


\begin{itemize}
  \item No authentication on manual control. Any device on the local Wi-Fi network can
issue control commands, with no login, no role check, and no protection against
unauthorised or accidental interference.
\end{itemize}

\begin{itemize}
  \item No occupancy verification beyond environmental sensors. The system cannot dis-
tinguish a scheduled-but-empty laboratory from an actually occupied one, since it
has no camera or vision-based confirmation.
\end{itemize}

\begin{itemize}
  \item No backend, persistence, or audit trail. The ESP32 operates in isolation; there
is no central record of historical sensor data, override events, or system state for
administrative review.
\end{itemize}



\subsubsection*{Relevance to GreenLab}

device itself to an authenticated backend, closing the most significant safety and governance gap this project leaves open. The sensor set is also extended with the HC-SR04 and
the YOLO-based camera pipeline to address the occupancy-verification limitation identified
above.





\textit{Figure 2.3: ESP32-based classroom automation hardware prototype reported by Ingole et al.}




\subsection{SmartHaus --- Angular and REST API IoT Dashboard}


Methodology and Technical Summary SmartHaus is an open-source home automation system for monitoring and controlling IoT devices through real-time dashboards and
a RESTful API [?]. Its repository describes an Angular client, a Node.js server, REST

27


communication, microcontroller compatibility with Arduino, ESP8266, and ESP32, and
API functions for reading and writing digital and analogue pins. The project is relevant to
GreenLab because GreenLab also relies on a dashboard-centred architecture in which users
monitor device state and issue authorised control commands. SmartHaus demonstrates that
an Angular single-page application can serve as an effective IoT monitoring and control
interface, with real-time device status updates and responsive command submission [?].

The main contribution of SmartHaus to GreenLab is architectural rather than domainspecific. It illustrates how a web dashboard can expose device-level operations through
REST abstractions while keeping IoT devices on a local network. The REST API decouples the dashboard from the device communication protocol, allowing the dashboard to be
developed and deployed independently of the microcontroller firmware. This decoupling is
essential for GreenLab because the Angular dashboard must work with the ASP.NET Core
backend rather than directly with the ESP32, and the backend must enforce business rules,
authentication, and audit logging that are not present in the SmartHaus architecture.

GreenLab follows a similar separation between dashboard controls and embedded
devices but adapts the pattern for faculty laboratories. The SmartHaus system treats all
users as equivalent controllers with access to all devices, which is appropriate for a home
environment but not for a faculty laboratory where instructors should only control their
assigned laboratories, administrators should manage schedules and approvals, and some
actions should require explicit authorisation and audit logging. GreenLab adds role-based
access that distinguishes instructor, administrator, and service roles; laboratory scheduling
and transfer approvals that prevent unauthorised room changes; audit logs that record every
state-changing action with actor, timestamp, and outcome; and YOLO occupancy integration that provides automated occupancy-aware control beyond simple on-off switching.



\subsubsection*{Advantages}


\begin{itemize}
  \item Validated Angular + REST pattern. Demonstrates a working, decoupled architec-
ture in which an Angular client communicates with IoT devices entirely through a
REST abstraction layer, directly informing GreenLab's frontend-to-backend design.
\end{itemize}

\begin{itemize}
  \item Broad microcontroller compatibility. Supports Arduino, ESP8266, and ESP32 de-
vices, showing that the REST abstraction generalises across common low-cost mi-
\end{itemize}

28
GreenLab: Intelligent Laboratory Energy Management System


crocontroller platforms.

\begin{itemize}
  \item Open-source transparency. As an open repository, the architecture and API design
can be inspected directly, giving a verifiable, citable technical reference rather than a
description relying solely on a paper's claims.
\end{itemize}



\subsubsection*{Disadvantages / Limitations}



\begin{itemize}
  \item No role differentiation. All dashboard users have equivalent access to all connected
devices, which is unsuitable for a multi-instructor, multi-administrator faculty envi-
ronment with differentiated authority.
\end{itemize}

\begin{itemize}
  \item No domain-specific automation logic. SmartHaus exposes raw pin read/write oper-
ations rather than higher-level automation concepts such as occupancy classification,
lighting decision matrices, or fan activation tables.
\end{itemize}

\begin{itemize}
  \item No audit logging or accountability mechanism. State-changing actions are not
recorded in a persistent, reviewable audit trail.
\end{itemize}



\subsubsection*{Relevance to GreenLab}

SmartHaus --- keeping the dashboard's REST contract independent of the embedded device
protocol --- but replaces SmartHaus's flat, single-tier user model with role-based access
control, laboratory-scoped permissions, transfer approval workflows, and comprehensive
audit logging, none of which the SmartHaus architecture provides.



\section{Comparative Analysis}


The reviewed sources collectively show that existing works address individual parts of
smart-building automation, including detection, control, or dashboards, but none combines
all the governance features that a real faculty laboratory requires. Chaudhari et al. provide
a survey of sensing technologies rather than a deployed system, and Zhang et al. contribute
an improved YOLO algorithm rather than a complete automation platform; consistent with
Sections 2.2.2 and 2.2.3, these two works are discussed in the text but excluded from the
comparison table below. The remaining six reviewed works and GreenLab are compared

29


across the technical dimensions that most directly distinguish GreenLab from the existing
body of work.




30
GreenLab: Intelligent Laboratory Energy Management System
Table 2.1: Comprehensive Feature Comparison Matrix
Dimension Challa et al. Waghchoure et Wibowo et al. Ahmed et al. Ingole et al. SmartHaus GreenLab
(2025) al. (2025) (2023) (2020) (2024)
Target environ- Classroom Classroom University labora- Classroom Classroom Generic home/IoT Faculty laboratory
ment tory
Sensing (multi- $\checkmark$ Camera + gas $\checkmark$ PIR + camera $\checkmark$ IoT sensor net- $\checkmark$ Motion + envi- $\checkmark$ LDR + tempera- × No environmen- $\checkmark$ LDR + DHT11 +
modal) sensor work ronmental ture tal sensing HC-SR04
AI / Computer vi- $\checkmark$ YOLOv4 (face- Partial (camera × × × × $\checkmark$ YOLO-based
sion based) verification, no person counting
count)
Automated light- × $\checkmark$ $\checkmark$ $\checkmark$ $\checkmark$ Partial (manual pin $\checkmark$ Decision-matrix
ing control write) driven
Automated venti- $\checkmark$ HVAC-linked $\checkmark$ Fans Partial $\checkmark$ Fans $\checkmark$ Fans Partial (manual pin $\checkmark$ Fan activation ta-
lation / HVAC con- write) ble
trol
Web dashboard × × $\checkmark$ Monitoring only $\checkmark$ Monitoring + Partial (Wi-Fi $\checkmark$ Angular dash- $\checkmark$ Angular SPA
basic override command only, no board dashboard
UI)
Role-based access × × × × × × $\checkmark$ Instructor / Ad-
control (Gover- ministrator / Service
nance) roles
Manual override × × Basic, no expiry Basic, no expiry $\checkmark$ Wi-Fi only, × $\checkmark$ Time-limited,
with expiry unauthenticated authenticated, au-
ditable
Audit logging / ac- × × × × × × $\checkmark$ Tamper-resistant
countability audit log
Laboratory trans- × × × × × × $\checkmark$ Request and ap-
fer workflow proval workflow
Backend persis- Partial (local pro- × $\checkmark$ Centralised Partial × $\checkmark$ Node.js server $\checkmark$ ASP.NET Core +
tence layer cessing only) monitoring SQL Server

Sources: Challa et al. [?]; Waghchoure et al. [?]; Wibowo et al. [?]; Ahmed et al. [?]; Ingole et al. [?]; SmartHaus [?].




31


The comparison table confirms that no single reviewed source combines laboratory
scheduling, role-based dashboard governance, ESP32 relay control, YOLO occupancy detection, manual override with audit trails, and accountability logging. Challa et al. provide
strong computer-vision occupancy detection but no web-based governance. Wibowo et al.
target university laboratories but lack YOLO integration and role-based access. Ahmed et
al. demonstrate classroom automation but without computer vision or audit logging. Ingole et al. provide the closest hardware match but operate an isolated ESP32 controller
without backend authority or governance. SmartHaus validates the dashboard-and-REST
architectural pattern but provides no domain-specific automation logic and no role differentiation. GreenLab fills this gap by integrating sensing, AI-based detection, automated
control, dashboard governance, and audit accountability within a single layered architecture designed specifically for faculty laboratory deployment.



\section{Literature Gap Addressed by GreenLab}


The literature review reveals a clear gap in the existing body of work: no single reviewed
source combines laboratory scheduling, role-based governance, ESP32 relay control, YOLO
occupancy detection, manual override with time-limited expiry and automatic policy restoration, and audit logging with tamper-resistant operational records. Each reviewed source
addresses one or two aspects of smart-building automation, but the governance layer required for a real faculty deployment is consistently absent. Occupancy detection papers
focus on algorithmic accuracy without addressing who may view or act on the detection
results. Classroom automation projects focus on sensor-to-relay control without providing
role-based access or audit trails. Dashboard projects focus on user interface and REST
architecture without integrating occupancy-aware automation or laboratory scheduling.

This gap can be stated precisely along three structural lines, each grounded in the eight
reviewed sources above:


\begin{itemize}
  \item The detection-without-governance gap. Challa et al. [?], Chaudhari et al. [?],
and Zhang et al. [?] demonstrate that occupancy can be detected accurately and
efficiently through multimodal sensing and lightweight computer vision, but none
of these sources extends detection into an accountable control system: there is no
\end{itemize}

32
GreenLab: Intelligent Laboratory Energy Management System


concept of who is authorised to act on a detection result, and no record of who did.

\begin{itemize}
  \item The automation-without-accountability gap. Waghchoure et al. [?], Ahmed et
al. [?], and Ingole et al. [?] demonstrate that sensor-to-relay automation is achiev-
able with low-cost embedded hardware, but each treats manual override as an unau-
thenticated, unaudited convenience feature rather than a governed operation, leaving
unauthorised or unaccountable control as an open risk in any production deployment.
\end{itemize}

\begin{itemize}
  \item The dashboard-without-domain-logic gap. Wibowo et al. [?] and SmartHaus [?]
demonstrate that centralised web dashboards can provide visibility and control over
distributed IoT devices, but neither integrates occupancy-aware automation logic,
laboratory-specific scheduling, or the role separation that a multi-instructor, multi-
administrator faculty environment requires.
\end{itemize}


GreenLab addresses this three-part gap by designing a system in which energy management is not merely a sensor-to-relay control problem but a governance problem that
requires authentication, authorisation, override management, transfer approvals, fault handling, and audit logging. The four-tier architecture ensures that each concern is handled by
the appropriate tier: the ESP32 handles local sensor reading and relay actuation, the Python
YOLO service handles occupancy inference, the ASP.NET Core backend handles business
rules and persistence, and the Angular dashboard handles presentation and user interaction.
This integration of automation and governance is the distinguishing contribution of GreenLab relative to the reviewed literature, and it is motivated by the practical requirements of
faculty laboratories where automation must coexist with teaching authority, administrative
oversight, and safe equipment operation.




33
